\documentclass[final]{beamer}
% you can chose your theme here:
\usetheme{I6dv}
% further beamerposter themes are available at 
% http://www-i6.informatik.rwth-aachen.de/~dreuw/latexbeamerposter.php
  \usepackage{type1cm}
  \usepackage{calc} 
  \usepackage{times}
  \usepackage{amsmath,amsthm, amssymb, latexsym}
\usepackage[utf8]{inputenc}  %or: \usepackage[latin1]{inputenc}
  \boldmath
  \usepackage[english]{babel}
  \usepackage[orientation=portrait,size=a1,scale=1.5,debug]{beamerposter}
\title{gLite sur Grid'5000:\\[0.1em] vers une plate-forme d'expérimentation\\[0.1em] à taille réelle pour les grilles de production}
\author{Sébastien Badia et Lucas Nussbaum}
\institute{LORIA / Nancy-Université \hskip 2em INRIA Nancy -- Grand Est}

%  \newcommand{\footlinetext}{Lehrstuhl f\"ur Informatik 6 - Computer Science Department - RWTH Aachen University - Aachen, Germany \par Mail: \texttt{<surname>@cs.rwth-aachen.de} \hfill WWW: \texttt{http://www-i6.informatik.rwth-aachen.de}\vskip1ex}
%  \date{Jul. 31th, 2007}

\setbeamerfont{itemize/enumerate subbody}{size=\normalsize}
\setbeamerfont{itemize/enumerate subsubbody}{size=\normalsize}

\setbeamertemplate{footline}{
  \begin{beamercolorbox}[wd=\paperwidth]{upper separation line foot}
    \rule{0pt}{2pt}
  \end{beamercolorbox}

  \begin{beamercolorbox}{footline}%
    \vskip 0.3em
    \begin{columns}[t]
	    \begin{column}{0.24\linewidth}
		    \center		    \includegraphics[height=3em]{grid5000}
	    \end{column}

	    \begin{column}{0.24\linewidth}
		    \center	    \includegraphics[height=3em]{logo_loria_complet_couleur-crop}
	    \end{column}

	    \begin{column}{0.24\linewidth}
		    \center    \includegraphics[height=3em]{NancyU}
	    \end{column}
	    \begin{column}{0.24\linewidth}
		    \center		    \includegraphics[height=3em]{INRIA_SCIENTIFIQUE_FR_CMJN_CROP.pdf}
	    \end{column}

    \end{columns}
    \vskip 1em
  \end{beamercolorbox}
}

  \begin{document}
  \begin{frame}{} 
	  \begin{center}
    \begin{columns}[t]
      \begin{column}{.48\linewidth}

	\begin{block}{Objectifs}
		\begin{itemize}
			\item \textbf{Utiliser Grid'5000 comme infrastructure de test pour gLite}
					\bigskip
			\begin{itemize}
				\item Pouvoir reproduire les conditions expérimentales nécessaires (exemple: simuler un crash de services)
					\bigskip
				\item Pouvoir remplacer des parties de l'infrastructure pour tester de nouvelles versions de composants du middleware

					\bigskip
				\item Pouvoir réaliser et reproduire des expériences dans un environnement stable
					\bigskip
				\item Éviter de perturber ou de surcharger l'infrastructure de production
			\end{itemize}
	\end{itemize}
	\end{block} \vskip 1em

	\begin{block}{Grid'5000}
		\begin{itemize}
			\item Plate-forme expérimentale pour la recherche sur les systèmes distribués et le calcul parallèle à hautes performances
					\bigskip
			\item 1700 machines (7000 coeurs) dans 10 sites en France
					\bigskip
			\item Reconfiguration: remplacement du système installé sur les noeuds, isolation via VLANs, etc.
		\end{itemize}
		\bigskip
		\centerline{\includegraphics[width=0.3\columnwidth]{g5k}}
	\end{block}  \vskip 1em

	\begin{block}{Infrastructure gLite déployée}
\begin{itemize}

	\item une \textbf{VO} et son \textbf{VOMS (\textsl{Virtual Organization Membership Service})},
	annuaire des utilisateurs ;
					\bigskip

\item plusieurs sites, composés de:

		\bigskip
\begin{enumerate}

	\item un \textbf{BDII (\textsl{Berkeley Database Information Index})}, annuaire
		des ressources disponibles sur chaque site ;

		\bigskip

	\item un \textbf{CE (\textsl{Computing Element})}, service de soumission des
		tâches à un site de calcul donné, et un ou plusieurs clusters composés de noeuds de calcul ;

		\bigskip
	\item un \textbf{WMS (\textsl{Workload Management System})}, méta-ordonnanceur entre les sites de la grille; dans le cadre de
		notre travail, le couple \textbf{Torque/Maui} a été utilisé ;
		\bigskip

	\item une \textbf{UI (\textsl{User Interface})}, interface d'accès pour les
		utilisateurs.

\end{enumerate}
\end{itemize}
	\end{block}

      \end{column}
      \begin{column}{.48\linewidth}
\begin{block}{Outils développés}
	\begin{itemize}
		\item \textbf{Image Scientific Linux 5.5} minimale générique (fonctionnant sur tous les clusters de Grid'5000) ;
		\bigskip
		\item \textbf{Scripts Ruby} permettant une installation automatique de gLite à partir des dépôts RPM
		\bigskip
		\begin{itemize}
			\item Génération d'une autorité de certification pour générer et signer automatiquement les certificats des utilisateurs et des machines
		\bigskip
			\item Pré-remplissage du cache RPM sur les noeuds avec BitTorrent
		\end{itemize}
	\end{itemize}
	\begin{center}
		\textbf{\url{https://github.com/sbadia/gdeploy/}}
	\end{center}
	\end{block} \vskip 1em


\begin{block}{Résultats}
	\begin{center}
		\textbf{Utilisation de Grid'5000 pour déployer le middleware gLite}
	\end{center}
\begin{itemize}
	\item Déploiement sur \textbf{5 sites, 8 clusters et 441 machines}
		\bigskip
\item Installation de l'ensemble des machines sous Scientific Linux 5.5 minimale avec Kadeploy: 10 minutes
		\bigskip
\item Configuration automatique de l'infrastructure gLite: \textbf{FIXME minutes}
\end{itemize}
	\end{block} \vskip 1em

	\begin{block}{Perspectives}
	\begin{itemize}
		\item Déploiement de plusieurs VO
		\bigskip
		\item Déploiement de services gLite supplémentaires
		\bigskip
			\begin{itemize}
				\item Stockage
		\bigskip
				\item Monitoring
		\bigskip
			\end{itemize}
		\item \textbf{Collaborations ?}
		\bigskip
			\begin{itemize}
				\item \textbf{Expériences sur des évolutions de composants du middleware gLite}
		\bigskip
				\item \textbf{Expériences sur des outils interagissant avec le middleware gLite: moteurs de workflows, gestionnaires de jobs pilotes, etc.}
		\bigskip
					\begin{itemize}
						\item Simulation de pannes de services
		\bigskip
						\item Injection de charge
		\bigskip
						\item Soumission d'un grand nombre de tâches factices
					\end{itemize}
			\end{itemize}
	\end{itemize}
	\end{block}
      
      \end{column}
    \end{columns}
 \vfill
    \end{center}
  \end{frame}
\end{document}


%%%%%%%%%%%%%%%%%%%%%%%%%%%%%%%%%%%%%%%%%%%%%%%%%%%%%%%%%%%%%%%%%%%%%%%%%%%%%%%%%%%%%%%%%%%%%%%%%%%%
%%% Local Variables: 
%%% mode: latex
%%% TeX-PDF-mode: t
